\scanpage{136}
\mspage{7}
% archive sheet number: 135

et un hom
\diagram{\begin{tikzcd}[ampersand replacement=\&,column sep=small]
A\&\overset{p}{\unsure{{-}|\longrightarrow}}\& A/N
\end{tikzcd}}
ayant le noyau $N$ (ou $\supset N$ ?), qui soit \uline{universel}, en ce sens que tout hom d'ateliers
\diagram{\begin{tikzcd}[ampersand replacement=\&,column sep=small]
A\&\overset{f}{\unsure{{-}|\longrightarrow}}\& A'
\end{tikzcd}}
ayant un noyau $\supset N$, se factorise par $p$ de façon unique. (\uline{NB} La notion de \uline{composition} est claire, en interprétant les hom. d'ateliers comme des applications ensemblistes
\diagram{\begin{tikzcd}[ampersand replacement=\&,column sep=small]
A\amalg\{\varnothing\}\&\overset{\widehat{f}}{\unsure{{-}|\longrightarrow}}\& A'\amalg\{\varnothing\}
\end{tikzcd}}
transformant $\{\varnothing\}$ en $\{\varnothing\}$, et \deleted{\unsure{un hom}} \deleted{\unsure{que}} dont la restriction à $A$ a les propriétés préc. \unsure{cela revient} à dire que $\widehat{f}$ est compatible à $\trianglelefteq,\overset{\circ}{\ll},\mathrel{|\circ|}$ sur $\widehat{A}$ !
% margin note written sideways
% slanted
\margin{\unsure{commencer} par ça !}

L'exemple type est celui où
\[
A=\unsure{\operatorname{Fig\acute{e}lem}|\operatorname{Fig\acute{e}l\acute{e}m}}(E),\quad N=\{X\in A\mid |X|\subset E\setminus E')
\]
\[
A/N=\unsure{\operatorname{Fig\,elem}}(E'),\quad\text{mais}
\]
\unsure{on veut} \unsure{trouver|toujours} \inserted{dans ce cas} \unsure{interprété|interpréter} comme quotient de $A\setminus N$
